\documentclass[11pt]{article}
\usepackage[margin=1in]{geometry}
\usepackage{microtype}
\usepackage{hyperref}
\usepackage{url}
\usepackage{amsmath}
\usepackage{amssymb}
\usepackage{graphicx}
\usepackage{booktabs}
\usepackage{xcolor}
\usepackage[numbers,sort&compress]{natbib}

\hypersetup{
  colorlinks=true,
  linkcolor=teal,
  urlcolor=teal,
  citecolor=teal,
}

\title{\texttt{ghostloop}: A Post-Hoc Analysis Layer for Embodied Agents}
\author{
  Joe Munene \\
  Complex Developers \\
  \texttt{joemunene984@gmail.com}
}
\date{May 2026}

\begin{document}
\maketitle

\begin{abstract}
LLM-agent frameworks have converged on a stable shape: a tool registry,
a policy pipeline, and a structured trace log. Robotics frameworks have
not. While ROS 2 provides middleware and Vision-Language-Action (VLA)
models provide policies, no widely-used layer sits between them with
the analytical surface that LLM-safety practice now expects. We
describe \texttt{ghostloop}, an open-source agent runtime for
embodied AI built around a fail-closed safety pipeline and a typed
trace, and contribute a \textbf{post-hoc analysis tier} that is
absent from contemporary robotics frameworks: counterfactual trace
replay, leave-one-out causal failure attribution, LLM-as-judge
trace scoring, declarative property mining, and adversarial Episode
search via CMA-ES. We position these as natural transplants from
LLM safety into embodied AI, formalise their semantics over a
\texttt{TraceEvent} DAG, and demonstrate them on six robot
morphologies (arm, mobile base, quadruped, drone, mobile-arm,
humanoid) through six backends (Mock, MuJoCo, PyBullet, Gymnasium,
ROS 2, RandomizedBackend). The runtime ships 333 tests across
eleven releases and is publicly available at
\url{https://github.com/joemunene-by/ghostloop} under the MIT
license.
\end{abstract}

\section{Introduction}

Robotics in 2026 has two healthy ecosystems with a missing middle.
ROS 2 \citep{macenski2022ros2} provides middleware: a message bus,
lifecycle management, drivers, navigation. VLA models such as
OpenVLA \citep{kim2024openvla}, RT-2 \citep{brohan2023rt2}, $\pi_0$
\citep{black2024pi0}, and Octo \citep{octo2024} provide policies:
vision-and-language conditioned action heads. Neither layer exposes
what LLM-agent practice now treats as standard: a typed tool
registry, a fail-closed policy pipeline, paired-comparison eval with
McNemar exact tests \citep{mcnemar1947note}, and post-hoc analytical
tooling for traces.

We propose \texttt{ghostloop}, an open-source runtime that fills this
gap. The runtime borrows shape from LLM-agent frameworks but binds
to robot primitives instead of API tools, with a backend abstraction
spanning sim and hardware. The contribution of this paper is not the
runtime itself: many such runtimes exist as ad-hoc internal tooling
in robotics labs. Our contribution is the \textbf{post-hoc analysis
tier} that comes with it, transplanted from LLM-safety practice and
made native to a robotics trace shape.

Specifically:
\begin{enumerate}
  \item \textbf{Counterfactual trace replay} (\S\ref{sec:counterfactual}):
        re-run a recorded trace's observations through a candidate
        policy and diff its decisions step-by-step.
  \item \textbf{Causal failure attribution} (\S\ref{sec:causal}):
        leave-one-out ablation that ranks events by necessity for
        a property violation.
  \item \textbf{LLM-as-judge trace scoring} (\S\ref{sec:judge}): score
        recorded episodes against a four-criterion rubric using any
        OpenAI-compatible chat client.
  \item \textbf{Declarative property mining} (\S\ref{sec:mining}):
        discover candidate invariants from a corpus of successful
        traces and promote them to formal \texttt{Property} objects.
  \item \textbf{Adversarial Episode search} (\S\ref{sec:adversarial}):
        random / grid / CMA-ES search over Episode initial states
        to find failure-prone seeds.
\end{enumerate}

Each pillar exists in adjacent fields (LLM red-teaming, formal
verification, mutation testing) but, to our knowledge, has not been
unified inside a robotics framework. We argue this unification is
the missing piece: as VLA models become the dominant policy
substrate, the runtime around them must adopt the analytical rigor
LLM-safety has been building for half a decade.

\section{Related Work}

\paragraph{Agent runtimes for LLMs.} The pattern of typed tool
registry plus policy pipeline plus structured trace appears in
LangChain \citep{langchain2022}, LlamaIndex \citep{llamaindex2023},
the AutoGen agent framework \citep{wu2023autogen}, and Anthropic's
Model Context Protocol (MCP) \citep{anthropic2024mcp}. Our
\texttt{Runtime} mirrors this shape exactly; the novelty is binding
it to robot \texttt{Primitive}s rather than HTTP / Python tools.

\paragraph{Safety in embodied AI.} Constraint-based safety appears
in OMPL motion planning \citep{sucan2012ompl}, formal-method
verification of controllers \citep{kress2017synthesis}, and
control-barrier-function approaches \citep{ames2019cbf}. None of
these sit at the agent-loop layer; they target the controller
or planner. Our \texttt{PolicyPipeline} layer is upstream of all
of them and gates the LLM's intent emission directly.

\paragraph{Counterfactual reasoning in ML.} Counterfactual
explanations \citep{wachter2017counterfactual} and counterfactual
reinforcement learning \citep{buesing2018woulda} are well studied
for decision-making models; \citet{anthropic2024sleeper}
demonstrate counterfactual prompting in LLM safety auditing. We
transplant the technique to a robot trace: given a trace from
policy $A$, ask what policy $B$ would have done at each step.

\paragraph{Causal attribution.} Shapley-value attribution for ML
\citep{lundberg2017shap} and counterfactual root-cause analysis
\citep{ji2023rca} are now standard for ML diagnostics. Robotics
incident analysis remains predominantly manual.

\paragraph{LLM-as-judge.} Originally validated for chatbot
preference modelling \citep{zheng2023judging,gu2024llmjudge}, the
LLM-as-judge methodology has spread to code-generation eval and
RAG quality. Applying it to robot trace eval against a structured
safety + task rubric is, to our knowledge, novel.

\paragraph{Adversarial testing in robotics.} Domain randomization
for sim-to-real \citep{tobin2017dr} randomises uniformly; our
adversarial Episode search maximises an explicit failure objective
via CMA-ES \citep{hansen2003cmaes}, making the test budget
adaptive.

\section{Runtime Architecture}

\subsection{Core abstractions}

A \texttt{Runtime} composes a \texttt{Backend} (sim or hardware
adapter), a \texttt{PrimitiveRegistry} (named callables), and a
\texttt{PolicyPipeline} (ordered fail-closed gates). One step
runs:

\begin{verbatim}
  intent     -> registry.get(intent.name) -> primitive
  primitive  -> pipeline.check(intent, primitive) -> Decision
  if ALLOW   -> primitive.call(backend, **intent.args) -> Result
  always     -> trace.append(TraceEvent(intent, decision, result, ...))
\end{verbatim}

A \texttt{TraceEvent} carries: \texttt{step}, \texttt{intent},
\texttt{decision}, \texttt{result}, \texttt{state\_before},
\texttt{state\_after}, and a \texttt{timestamp}. The
\texttt{Trace} is append-only and JSON-serialisable. This typed
event shape is the substrate for every analytical operator in
\S\ref{sec:posthoc}.

\subsection{Backends and morphologies}

Six backends ship: \texttt{MockBackend} (in-memory, no install),
\texttt{MuJoCoBackend} \citep{todorov2012mujoco},
\texttt{PyBulletBackend} \citep{coumans2016pybullet},
\texttt{GymnasiumBackend} \citep{towers2023gymnasium},
\texttt{ROS2Backend}, and a \texttt{RandomizedBackend} wrapper that
applies position noise, timing jitter, action dropout, and
sticky-action repeats for sim-to-real \citep{tobin2017dr}.

A cross-morphology primitive library covers seven categories
(\texttt{MOBILE\_BASE}, \texttt{QUADRUPED}, \texttt{HUMANOID},
\texttt{AERIAL}, \texttt{DEXTEROUS}, \texttt{SENSING},
\texttt{GENERIC}) so a single \texttt{RobotProfile} declaration
adapts the runtime from a Franka 7-DOF arm to a Spot quadruped or a
Tello quadcopter without code changes.

\subsection{Safety pipeline}

Twelve gates ship: \texttt{DenyListGate}, \texttt{RateLimitGate},
\texttt{GeofenceGate}, \texttt{ForceCapGate}, \texttt{Cooldown\-Gate},
\texttt{TimeWindowGate}, \texttt{ActionSmoothingGate} (velocity /
acceleration limits), \texttt{HumanInTheLoopGate}, \texttt{Obstacle\-Avoid\-anceGate},
\texttt{RetryPolicy}, plus \texttt{LLMPolicy} and \texttt{VLAPolicy}
adapters. Gates compose; the pipeline short-circuits on the first
\texttt{DENY}.

\section{Post-Hoc Analysis Tier}
\label{sec:posthoc}

This is the contribution. The \texttt{TraceEvent} substrate makes
each operator below a small typed function over a \texttt{Trace}
object.

\subsection{Counterfactual trace replay}
\label{sec:counterfactual}

Given a recorded \texttt{Trace} $T$ from policy $\pi_A$ and a
candidate policy $\pi_B$, we walk every event of $T$ and ask
$\pi_B$ what intent it would emit given the same
\texttt{state\_before}. The output is a \texttt{Counterfactual\-Trace}:

\[
  \text{CF}(T, \pi_B) = \{(\text{state}_i, \pi_A(\cdot)_i, \pi_B(\text{state}_i)) : i \in T\}
\]

with metrics \emph{divergence rate} (fraction of events where
$\pi_A$ and $\pi_B$ disagree) and \emph{first-divergence step}
(earliest disagreement). We do not re-execute the world; this is
shadow-mode reasoning, the same shape used in LLM safety auditing
\citep{anthropic2024sleeper}. Re-execution against a sim backend
is straightforward via the standard \texttt{Runtime} loop.

\subsection{Causal failure attribution}
\label{sec:causal}

Given a failing trace $T$ for property $\phi$, for each event
$e_i \in T$ we evaluate $\phi(T \setminus \{e_i\})$. Each
$e_i$ receives a \emph{necessity} score:

\[
  \text{necessity}(e_i) = \begin{cases}
    1.0 & \text{if } \phi(T \setminus \{e_i\}) \text{ holds} \\
    0.5 \cdot \frac{|V_T| - |V_{T \setminus \{e_i\}}|}{|V_T|} & \text{otherwise}
  \end{cases}
\]

where $V_T$ is the set of violations $\phi$ records on $T$. Events
whose removal restores $\phi$ are causally necessary; events whose
removal merely reduces violation count receive partial credit. A
top-$k$ ranking surfaces likely root causes. We additionally
provide \texttt{minimal\_cause\_set} for greedy multi-event
attribution.

\subsection{LLM-as-judge}
\label{sec:judge}

Following \citet{zheng2023judging}, we summarise a trace and pass
it to a chat client with a four-criterion rubric:
\textit{task\_completion}, \textit{safety\_adherence},
\textit{efficiency}, \textit{recoverability}, each scored in
$[0,1]$. The judge returns structured JSON; the wrapper validates
and exposes the sub-scores plus a single composite. A
\texttt{HeuristicJudge} fallback ships for air-gapped CI: weighted
predicate scoring (no-violations, fraction-allowed, fraction-OK,
step-count) gives a scoring shape compatible with the LLM judge
without an external model call.

\subsection{Declarative property mining}
\label{sec:mining}

Given a corpus of successful traces $\mathcal{T} = \{T_1, \dots, T_n\}$,
three pattern miners run:

\paragraph{Followup transitions.} For every pair of primitive
names $(x, y)$, count traces where $x$ is followed by $y$ within
window $W$. If support and confidence both exceed threshold $\tau$,
emit a \texttt{MinedProperty} that fails when $x$ is not followed
by $y$ within $W$ on a future trace.

\paragraph{Numeric bounds.} For each named observation field,
take the empirical max across $\mathcal{T}$, multiply by margin
factor $\mu$, and emit a never-exceeds property at that threshold.

\paragraph{Workspace bounds.} For every recorded position, take
the empirical bounding box across $\mathcal{T}$ and emit a
\texttt{StaysInsideWorkspace} candidate.

Each \texttt{MinedProperty} carries support, confidence, and a
\texttt{promote()} method returning a real \texttt{Property} the
operator can ratchet into the formal property engine.

\subsection{Adversarial Episode search}
\label{sec:adversarial}

Given an Episode template and a parameter range, three search
strategies emit failure-prone seeds:

\begin{itemize}
  \item \textbf{Random}: uniform sampling, baseline.
  \item \textbf{Grid}: exhaustive cartesian product.
  \item \textbf{CMA-ES}: \citep{hansen2003cmaes} biases toward
        regions with high failure score, where failure score $=$
        violation\_count $+ \mathbb{1}[\text{predicate failed}]$.
\end{itemize}

Worst-scoring seeds promote to the regression bench: each is now
a deterministic Episode the operator can re-run on every commit.

\section{Implementation}

\texttt{ghostloop} is implemented in pure Python 3.10+ with zero
runtime dependencies in the core. Optional extras lazy-import
backend libraries: \texttt{ghostloop[mujoco]},
\texttt{ghostloop[pybullet]}, \texttt{ghostloop[gym]},
\texttt{ghostloop[mcp]}, \texttt{ghostloop[dashboard]},
\texttt{ghostloop[otel]}. Conditional imports raise
\texttt{ImportError} with install hints at construction time only,
never at package import.

The codebase is 11 releases, 70+ modules, $\sim$7,000 lines of
Python, and 333 tests (8 live-gated for hardware libraries). Tests
cover every analytical operator above plus per-backend smoke
tests, a YAML-loadable robot profile system, and the MCP server
across stdio, streamable-HTTP, and SSE transports.

\section{Discussion and Limitations}

\paragraph{No re-execution counterfactuals yet.} The shadow-mode
counterfactual replay tells you what policy $\pi_B$ would have
\emph{chosen}; it does not tell you what would have happened. A
re-execution variant against a deterministic sim backend is
straightforward but adds computational cost; we leave it as a
future addition.

\paragraph{LLM-as-judge calibration.} We do not evaluate the
judge against human raters; the rubric was designed for clarity,
not validated against ground truth. Adopting the rater-pair
calibration procedure of \citet{zheng2023judging} on a labelled
trace corpus is the right next step.

\paragraph{Mining is statistical, not formal.} The mined
properties capture empirical regularities from a successful-trace
corpus. They will not generalise to behaviour distributions
different from the corpus. We treat them as \emph{candidate}
invariants that an operator promotes, not as automatic safety
guarantees.

\paragraph{No published numbers vs OpenVLA / $\pi_0$ here.}
\texttt{ghostloop} ships a \texttt{VLABenchmarkSuite} pre-populated
with published baselines from these papers, with the bench
harness running. We did not run a real OpenVLA-7B checkpoint
against the bench in time for this submission; reproducing those
numbers under the safety pipeline is the obvious follow-up.

\section{Conclusion}

We introduced \texttt{ghostloop}, an open-source agent runtime
for embodied AI that contributes a post-hoc analysis tier
transplanted from LLM-safety practice: counterfactual trace
replay, causal failure attribution, LLM-as-judge, property
mining, and adversarial Episode search. We argue this analytical
surface — already standard in LLM agent frameworks — is the
missing piece for robotics frameworks as VLA models become the
dominant policy substrate. The library is MIT-licensed at
\url{https://github.com/joemunene-by/ghostloop} and installable
via \texttt{pip install ghostloop}.

\paragraph{Acknowledgements.} This work began as a sibling
project to \texttt{GhostLM}, an 81M-parameter cybersecurity
language model trained from scratch
(\url{https://github.com/joemunene-by/GhostLM}). The agent-loop
shape and statistical bench harness are direct ports.

\bibliographystyle{plainnat}
\bibliography{refs}

\end{document}
